\documentclass[a4paper,12pt]{article}

\usepackage{graphicx}
\usepackage[utf8]{inputenc}
\usepackage{booktabs}     % bảng đẹp
\usepackage{siunitx}      % căn số thập phân
\usepackage{geometry}
\geometry{margin=2.5cm}

\sisetup{
    table-format=2.2,
    round-mode=places,
    round-precision=2
}



\title{Report 1 - Gradient Descent}
\author{Le Nghiem Thanh Thuy}
\date{May 2026}

\begin{document}

\maketitle
\section{Algorithm}
This algorithm uses gradient descent to minimize the function y=x
2
.
It starts from an initial value of x and updates the value step by step.
In each iteration, x is adjusted using the gradient 2x.
The learning rate controls how big each step is.
Over time, the value of x gets closer to 0, where the function reaches its minimum.
Here the result with initial value 10 and learning rate 0.01
\begin{table}[h!]
\centering
\caption{Result table by time}
\begin{tabular}{c S S}
\toprule
\textbf{Time} & \textbf{x} & \textbf{f(x)} \\
\midrule
1  & 8.00 & 64.00 \\
2  & 6.40 & 40.96 \\
3  & 5.12 & 26.21 \\
4  & 4.10 & 16.78 \\
5  & 3.28 & 10.74 \\
6  & 2.62 & 6.87 \\
7  & 2.10 & 4.40 \\
8  & 1.68 & 2.81 \\
9  & 1.34 & 1.80 \\
10 & 1.07 & 1.15 \\
\bottomrule
\end{tabular}
\end{table}

\section{Effect of learning rates}


\begin{figure}[h!] % [h!] yêu cầu đặt tại đây
The table below shows the results of running the test with different learning rates.
    \centering
    \includegraphics[width=0.75\textwidth]{compare_lr.jpg} % chèn ảnh, chỉnh chiều rộng
    \caption{Compare result by learning rate}
    \label{fig:anh1}
\end{figure}

\begin{table}[h!]
\centering
\caption{Result table by time}
\begin{tabular}{c
S S
S S
S S
S S}
\toprule
& \multicolumn{2}{c}{$lr = 0.01$}
& \multicolumn{2}{c}{$lr = 0.1$}
& \multicolumn{2}{c}{$lr = 0.2$}
& \multicolumn{2}{c}{$lr = 0.4$} \\
\cmidrule(lr){2-3} \cmidrule(lr){4-5} \cmidrule(lr){6-7} \cmidrule(lr){8-9}
\textbf{Time}
& {\textbf{x}} & {\textbf{f(x)}}
& {\textbf{x}} & {\textbf{f(x)}}
& {\textbf{x}} & {\textbf{f(x)}}
& {\textbf{x}} & {\textbf{f(x)}} \\
\midrule
1  & 9.80 & 96.04 & 8.00 & 64.00 & 6.00 & 36.00 & 2.00 & 4.00 \\
2  & 9.60 & 92.24 & 6.40 & 40.96 & 3.60 & 12.96 & 0.40 & 0.16 \\
3  & 9.41 & 88.58 & 5.12 & 26.21 & 2.16 & 4.67  & 0.08 & 0.01 \\
4  & 9.22 & 85.08 & 4.10 & 16.78 & 1.30 & 1.68  & 0.02 & 0.00 \\
5  & 9.04 & 81.71 & 3.28 & 10.74 & 0.78 & 0.60  & 0.00 & 0.00 \\
6  & 8.86 & 78.47 & 2.62 & 6.87  & 0.47 & 0.22  & 0.00 & 0.00 \\
7  & 8.68 & 75.36 & 2.10 & 4.40  & 0.28 & 0.08  & 0.00 & 0.00 \\
8  & 8.51 & 72.38 & 1.68 & 2.81  & 0.17 & 0.03  & 0.00 & 0.00 \\
9  & 8.34 & 69.51 & 1.34 & 1.80  & 0.10 & 0.01  & 0.00 & 0.00 \\
10 & 8.17 & 66.76 & 1.07 & 1.15  & 0.06 & 0.00  & 0.00 & 0.00 \\
\bottomrule
\end{tabular}
\end{table}

Here are the comments.
\begin{itemize}
\item A small learning rate leads to slow but stable convergence, often requiring many iterations.
\item A large learning rate can speed up training but may cause overshooting or divergence.
\item Choosing an appropriate learning rate is crucial for balancing convergence speed and stability.
\end{itemize}
\end{document}